% Conclusiones y recomendaciones de \nombreTercero (rúbrica del Entregable 3: 4 puntos, «del alumno»).
% Escribirlas en primera persona, con tus palabras: 2 o 3 párrafos alcanzan. Reemplaza la línea
% \pendiente{...} de abajo por tu texto. Preguntas guía, no hace falta responderlas todas:
%   - ¿Qué aprendiste sobre concurrencia con este trabajo (worker pool, barrera, Spin, speedup)?
%   - ¿Qué resultado te sorprendió y por qué?
%   - ¿Qué limitación ves en el trabajo y qué recomendarías para mejorarlo o escalarlo?
%   - En lo que hiciste tú (tu PR), ¿qué harías distinto?
% Para compilar y ver el PDF: cd tp/informe/tp && ./compilar.sh  (sale en build/main.pdf).
\subsubsection*{\nombreTercero}
Al probar casos borde de K-means aprendí que la concurrencia exige más que repartir viajes entre
goroutines. Cada punto debe procesarse una vez, los resultados parciales deben sumarse en un orden
fijo y los centroides solo pueden actualizarse después de la barrera. Las pruebas que aporté, con
más workers que puntos y con clusters vacíos, me ayudaron a comprobar que el algoritmo también
funciona fuera del caso habitual. Además, entendí por qué conviene combinar esas pruebas con Spin:
un resultado correcto en una ejecución no garantiza que todos los posibles órdenes de trabajo
respeten la sincronización.

Me sorprendió que el speedup no creciera al mismo ritmo que la cantidad de workers: con el dataset
completo fue de 3,80 con cuatro, 5,82 con ocho y 6,25 con dieciséis. Con 10.000 viajes, en cambio,
la versión concurrente fue más lenta que la secuencial. Concluyo que el costo de coordinación solo
se justifica cuando hay suficiente trabajo.

Recomiendo ampliar las pruebas con distintas relaciones entre workers y bloques. Así podríamos comprobar que el reparto sigue siendo correcto en
configuraciones que las pruebas actuales no cubren.
